% ===================================================================
%  TABEL Nr. 8 — Lõikus.
%  Traduction 2026 d'apres texto/io, controlee sur texto/fr et
%  texto/fi. Consigne : voir texto/et/10-tabel-01.tex.
%
%  LA SECONDE SCENE EST UNE VENDANGE, ET L'ESTONIE N'A PAS DE VIGNE.
%  C'EST LA MEILLEURE PAGE DE LA SERIE POUR VOIR CE QU'UNE LANGUE
%  FAIT D'UNE CHOSE QU'ELLE N'A RENCONTREE QUE PAR SON PRODUIT.
%
%  Trois boissons sur cette page, trois ages, trois formes de mot :
%    « õlu »       la biere, mot du fonds finno-ougrien, parce qu'on
%                  l'a toujours brassee a la maison ;
%    « vein »      le vin, emprunt a l'allemand, parce qu'il est
%                  arrive en tonneau ;
%    « viinamari » le raisin, qui est un compose : « viin »
%                  l'eau-de-vie et « mari » la baie — LA BAIE DONT ON
%                  FAIT DE L'ALCOOL.
%
%  ON NOMME LA PLANTE PAR CE QU'ON EN TIRE PARCE QU'ON N'A JAMAIS VU
%  QUE CE QU'ON EN TIRE. Un peuple qui n'a pas de vigne ne rencontre
%  pas la grappe au cep, il la rencontre dans une bouteille, et il
%  remonte du produit a la matiere.
%
%  ET C'EST LA FORME JUSTE D'UNE REGLE QUE LA SERIE CHERCHAIT DEPUIS
%  TROIS COLONNES. Le tableau 15 lituanien disait : la date d'arrivee
%  decide de la forme du nom. Le tableau 15 luxembourgeois corrigeait :
%  la date fixe qu'il y aura un mot neuf, pas lequel. Le tableau 15
%  romanche ajoutait : quand on emprunte, c'est la route de la chose
%  qui choisit le preteur. L'ESTONIEN DONNE LE TROISIEME CAS, CELUI
%  OU L'ON N'EMPRUNTE PAS ET OU L'ON NE COMPOSE PAS NON PLUS SUR LA
%  CHOSE : ON COMPOSE SUR SON USAGE. Non pas la baie-du-vin, non pas
%  la baie-du-sud, mais la baie-a-faire-boire.
%
%  ET LA NOTE (*) DE GUIGNON EST TOUJOURS LA, SUR CETTE PAGE-CI. Il y
%  manque un mot general pour « recolte » et le demande a
%  l'Academie. L'estonien ne le lui donnera pas : il a « lõikus » ce
%  qu'on coupe, « heinategu » ce qu'on fane, « korje » ce qu'on
%  cueille, et il en a fait ses noms de mois — « lõikuskuu » aout,
%  « heinakuu » juillet. Le tableau 3 l'avait deja montre ; ici c'est
%  l'auteur lui-meme qui pose la question a laquelle la langue avait
%  repondu avant qu'il l'ecrive.
%
%  ECART DECLARE DU BLOC t08-c2-03-1 : le fac-simile ido y ecrit
%  « gobleto 13) » sans la parenthese ouvrante. La coquille est du
%  compositeur de l'ido : la parenthese est retablie.
% ===================================================================

%%K t08-noto-1 noto
\VUnotes{143.2mm}{%
\textit{(*) Kuna see sõna ei ole täpne: kas on tegemist linalõikuse,
viinamarjakorje või õunakorjega? Ehk oleks hea kasutada sõna «} \VUgras{mesono}
\textit{», mida pakuti ajakirjas Mondo XIV, lk. 242. Kuid seni ei ole Akadeemia seda
uut tüve vastu võtnud ja see ootab arutelusid idistide tavapäraste reeglite
kohaselt.}%
}

%%K t08-apar-1 apar
\VUsaut{5.0mm}
\VUtitre{15.0pt}{60}{TABEL Nr. 8}
\VUsaut{3.0mm}
\VUcentre{13.2pt}{90}{\textit{Esimene stseen.}}
\VUsaut{3.0mm}
\VUpk{10.2pt}{40}{Lõikus (*).}
\VUsaut{4.0mm}
\VUpk{10.2pt}{40}{Maastiku ilmed.}
\VUsaut{3.0mm}

%%K t08-c1-01-1 p
1. --- \VUgras{Suvega} on maastiku ilme jälle muutunud. Vaole usaldatud seeme on
idanenud; roheline taimeke on mullast välja tulnud ja kasvatanud pea. Tera on
moodustunud, siis on see valminud, ja nüüdsest on vaja ainult lõigata.

%%K t08-c1-02-1 p
2. --- \VUgras{Põllulilled} on avanenud. \VUgras{Rukkililled} \textsuperscript{(1)},
\VUgras{moonid} \textsuperscript{(2)} ja \VUgras{sõrmkübarad} \textsuperscript{(3)}
segavad nisupõldude ühetaolisse tooni oma värskete \VUgras{õiekroonide} rõõmsat
vastandust.

%%K t08-c1-03-1 p
3. --- Viljapuud on end lehtedega ehtinud; vili on valminud. Väike \VUgras{kirsipuu}
\textsuperscript{(4)} on täis ilusaid \VUgras{kirsse} \textsuperscript{(5)}, mida lapsed
suure rõõmuga korjavad ja söövad.

%%K t08-c1-04-1 p
4. --- Nüüdsest võib kari veeta terve päeva vabas õhus.

%%K t08-c1-04-2 p
\VUgras{Talled} \textsuperscript{(6)} on suureks kasvanud ja neist on saanud ilusad
\VUgras{utid} \textsuperscript{(7)} ja ilusad \VUgras{oinad} \textsuperscript{(8)} tiheda
villaga. Nad karjatavad rahulikult \VUgras{karjamaa} \textsuperscript{(9)}
\VUgras{rohtu}, mis on kirju \VUgras{härjasilmadest}.

%%K t08-c1-04-3 p
Varsti hakatakse \VUgras{neid} loomi pügama, et saada nende \VUgras{villa}, millest
tehakse meie riiete kalevit.

%%K t08-tit-2 sub
\VUsaut{5.0mm}
\VUpk{10.2pt}{40}{Karjus.}
\VUsaut{3.4mm}

%%K t08-c1-05-1 p
5. --- \VUgras{Karjus} \textsuperscript{(10)} ei näi neid eriti tähele panevat. Ta
hoolib ainult äikesest, mis silmapiiril möödub, ja \VUgras{lambavaht}
\textsuperscript{(13)}, kes viib \VUgras{põlvpudeli} \textsuperscript{(14)} oma huultele,
näib veelgi muretum.

%%K t08-c1-06-1 p
6. --- Kuumus on rõhuv; sest ka \VUgras{koer} \textsuperscript{(15)} tuleb janu
kustutama; ta lakub \VUgras{oja} \textsuperscript{(16)} värsket vett ja hirmutab vaest
\VUgras{konna} \textsuperscript{(17)}, kes oli kaldarohtu peitunud. See hüppab selgesse
vette, kus õitsevad \VUgras{vesiroosid} \textsuperscript{(18)}.

%%K t08-c1-07-1 p
7. --- \VUgras{Västrik} \textsuperscript{(19)} suples \VUgras{allika}
\textsuperscript{(20)} juures, mis pritsib kahe kalju vahelt, ja peitub
\VUgras{okkaliste põõsaste} \textsuperscript{(21)} ja \VUgras{viirpuupõõsaste}
\textsuperscript{(22)} alla.

%%K t08-tit-3 sub
\VUsaut{5.0mm}
\VUpk{10.2pt}{40}{Lõikus.}
\VUsaut{3.4mm}

%%K t08-c1-08-1 p
8. --- Maalastele on nisulõikus väga lõbus aeg. Nad
\VUgras{kukerpallitavad}\textsuperscript{(24)} rõõmsalt \VUgras{õlekuhjadel}
\textsuperscript{(23)}, nad aitavad \VUgras{lõikajaid} \textsuperscript{(26)} ja sellal
kui need niidavad \VUgras{nisupõldu} \textsuperscript{(27)} oma \VUgras{vikatitega}
\textsuperscript{(25)}, mis on hästi \VUgras{luisuga} \textsuperscript{(30)} teritatud,
korjavad nemad nisu ja seovad selle \VUgras{vihkudeks} \textsuperscript{(31)} või
\VUgras{kimpudeks} \textsuperscript{(32)}.

%%K t08-c1-09-1 p
9. --- Mõnikord lubatakse suurematel neist \VUgras{sirbiga} \textsuperscript{(12)}
lõigata \VUgras{kuldseid kõrsi} \textsuperscript{(28)}, mis kannavad raskeid teradest
täis \VUgras{viljapäid} \textsuperscript{(29)}; ja väiksemad, valvates \VUgras{karja}
\textsuperscript{(33)}, mis karjatab \VUgras{karjamaal} \textsuperscript{(34)} suure
\VUgras{pärna} \textsuperscript{(35)} varjus, püüavad oma \VUgras{võrguga}
\textsuperscript{(36)} ilusaid \VUgras{liblikaid}.

%%K t08-c1-09-2 p
Mõned ronivad isegi \VUgras{vankrile} \textsuperscript{(37)}, et seada vihke, mida neile
ulatatakse \VUgras{hanguga} \textsuperscript{(38)}.

%%K t08-c1-10-1 p
10. --- Kogu \VUgras{tasandikul} \textsuperscript{(39)}, kus lendlevad \VUgras{lõokesed}
\textsuperscript{(41)}, valitseb elevus. Kõik on ametis lõikusega. Kuid sellal kui
vaesed jätkavad vanade riistade kasutamist, \VUgras{vikatite, sirpide} ja
\VUgras{kootide} \textsuperscript{(40)}, lasevad rikkad maaomanikud oma nisu või oma
\VUgras{rukki} \textsuperscript{(43)} niita \VUgras{niidumasinaga} \textsuperscript{(42)}.
Siis, ilma et oleks vaja vihke küüni viia, tehakse suured \VUgras{vihukuhjad}
\textsuperscript{(44)}.

%%K t08-c1-11-1 p
11. --- Siis tuuakse kohale \VUgras{rehepeksumasin} \textsuperscript{(47)}, mida
\VUgras{lokomobiil} \textsuperscript{(45)} paneb liikuma \VUgras{ülekanderihmaga}
\textsuperscript{(46)}, ja nii eraldatakse tera \VUgras{õlgedest}, ilma et lahkutaks
põllult, kuhu see külvati.

%%K t08-tit-4 sub
\VUsaut{5.0mm}
\VUpk{10.2pt}{40}{Torm.}
\VUsaut{3.4mm}

%%K t08-c1-12-1 p
12. --- Sageli tulevad lõikuse ajal \VUgras{tormid} \textsuperscript{(53)}. Kõigepealt
kuuldakse kaugelt \VUgras{kõue} mürinat; \VUgras{läänetorm} \textsuperscript{(54)} hakkab
puhuma ja painutab lõõtsudes \VUgras{metsa} \textsuperscript{(49)} kõige jämedamaid puid.
Mustad \VUgras{pilverünkad} \textsuperscript{(56)} tungivad taevasse; neid läbistavad
\VUgras{välgu} \textsuperscript{(55)} pimestavad siksakid. \VUgras{Vihm} ja mõnikord
\VUgras{rahe} hakkavad sadama, litsudes maha lõikuseks valmis vilja.

%%K t08-c1-13-1 p
13. --- Sageli süütab välk mõne ehitise. Nii põletati eelmisel aastal tuhaks
\VUgras{tuuliku} \textsuperscript{(50)} katus ja purustati kaks tema \VUgras{tiiba}
\textsuperscript{(51)}.

%%K t08-c1-14-1 p
14. --- Mõne tunni pärast sünnib rahu uuesti. Särav \VUgras{vikerkaar}
\textsuperscript{(52)} ilmub silmapiirile ja loodus võtab tagasi oma elu, mis oli mõneks
hetkeks katkenud. Maainimesed, kes olid varjunud \VUgras{onnidesse}
\textsuperscript{(48)}, hakkavad uuesti tööle ja püüavad julgelt parandada kahju, mida
nad äsja said.

%%K t08-tit-5 sub
\VUsaut{6.0mm}
\VUcentre{13.2pt}{90}{\textit{Teine stseen.}}
\VUsaut{3.6mm}

%%K t08-tit-6 sub
\VUpk{10.2pt}{40}{Jaht. --- Viinamarjakorje.}
\VUsaut{1.4mm}
\VUpk{10.2pt}{40}{Kalapüük.}
\VUsaut{4.0mm}

%%K t08-c2-01-1 p
1. --- Augusti lõpu või \VUgras{septembri} keskpaiga paiku, olenevalt piirkonnast, algab
jahihooaeg. Vaevalt on esimesed päikesekiired pannud \VUgras{sisalikud}
\textsuperscript{(2)} oma urgudest välja tulema, kui \VUgras{kütt} \textsuperscript{(5)}
asub teele oma \VUgras{koertega} \textsuperscript{(4)}.

%%K t08-c2-02-1 p
2. --- Ta on selga pannud oma tugeva paksust kalevist \VUgras{jahiülikonna}
\textsuperscript{(9)}, jalga nahksed \VUgras{säärised}, millega tal õnnestub käia
niiskes rohus, kuhu peituvad \VUgras{kärnkonnad} \textsuperscript{(1)}; ta on võtnud oma
\VUgras{püssi} \textsuperscript{(6)} ja oma \VUgras{jahikoti} \textsuperscript{(10)}, ja
vaat et ta läheb.

%%K t08-c2-03-1 p
3. --- Tagaajamise hoog viib ta keset viinamäge, kus viinamarjakorje on väga elav.
Viinamäekasvataja lapsed, kes tavaliselt valvavad teel kitsi, lasevad neil täna
juhuslikult karjatada, ja olles posti külge sidunud vana \VUgras{siku}
\textsuperscript{(3)}, kes ei jätaks kasutamata oma vabadust põgenemiseks, annavad nad
endale osa lõbust. Üks haarab viinamarjakobara \VUgras{korjekorvist}
\textsuperscript{(12)} ja lõbustab end, lastes \VUgras{mahla} \textsuperscript{(14)}
voolata \VUgras{peekrisse} \textsuperscript{(13)}, et teha virret (kääritamata veini).
Teine paneb pähe \VUgras{lehepärja} \textsuperscript{(15)}, mille ta äsja punus.

Nende lähedal tulevad häbematud \VUgras{varblased} \textsuperscript{(16)} isegi pressi
ette viinamarju röövima.

%%K t08-c2-04-1 p
4. --- Sellal kui lapsed lõbutsevad, töötavad mehed. \VUgras{Korjajad}
\textsuperscript{(18)} viivad viinamarjad keldrisse \VUgras{seljakorvidega}
\textsuperscript{(17)}; \VUgras{tündrisepp} \textsuperscript{(19)} parandab \VUgras{vaate}
\textsuperscript{(20)}, kontrollib, kas \VUgras{lauad} \textsuperscript{(22)} ja
\VUgras{põhjad} \textsuperscript{(21)} on korras, ja asendab katkised \VUgras{vitsad}
\textsuperscript{(23)}. Ta on ka teinud \VUgras{tõrred} \textsuperscript{(25)}, millesse
valatakse \VUgras{viinamarjad} \textsuperscript{(26)}, et neid kääritada.

%%K t08-c2-05-1 p
5. --- Kui kääritamine on lõppenud, lastakse vein välja ja jäägid pannakse
\VUgras{pressile} \textsuperscript{(28)}, ja järk-järgult \VUgras{kruviga}
\textsuperscript{(29)} pigistades pressitakse välja mahl, mis voolab \VUgras{vanni}
\textsuperscript{(32)}, ja nii saadakse veel \VUgras{veini} \textsuperscript{(31)}.

%%K t08-tit-8 sub
\VUsaut{6.0mm}
\VUpk{10.2pt}{40}{Õlu ja siider.}
\VUsaut{3.4mm}

%%K t08-c2-06-1 p
6. --- \VUgras{Keldri} \textsuperscript{(27)} müüril ronib \VUgras{humala}
\textsuperscript{(30)} väät. Siin on see ainult ilutaim, kuid mõnes maas, kus viinapuu
on väga haruldane, kasvatatakse humalat \VUgras{õlle} tegemiseks.

%%K t08-c2-07-1 p
7. --- Väike \VUgras{õunapuu} \textsuperscript{(33)} on koormatud õuntega, mis annavad,
kui nad on valminud, suurepärast \VUgras{siidrit}. Oma \VUgras{puuris}
\textsuperscript{(34)} vaatavad \VUgras{kanaarilind} \textsuperscript{(35)} ja
\VUgras{ohakalind} \textsuperscript{(36)} kadedate silmadega varblasi, kes vabalt
hullavad.

%%K t08-c2-08-1 p
8. --- Silmapiirini, küngaste nõlval, on ritta seatud \VUgras{viinapuu-teibad}
\textsuperscript{(40)}, mis toetavad imeilusaid \VUgras{viinapuu-tüvesid}
\textsuperscript{(39)}, ehitud raskete \VUgras{kobaratega}.

%%K t08-c2-08-2 p
Terve aasta on \VUgras{viinamäekasvataja} \textsuperscript{(42)} töötanud, et oma
\VUgras{viinamäge} \textsuperscript{(38)} hooldada, ja nüüd tasutakse talle vaeva eest
ilusa saagiga. \VUgras{Korjajannad} \textsuperscript{(41)} täidavad tõrred, mis on pandud
\VUgras{muulade} \textsuperscript{(43)} veetavale vankrile, ja kõik on rõõmsad.

%%K t08-tit-9 sub
\VUsaut{5.0mm}
\VUpk{10.2pt}{40}{Kalapüük.}
\VUsaut{3.4mm}

%%K t08-c2-09-1 p
9. --- Samamoodi on \VUgras{õngitsejatega} \textsuperscript{(44)}, kes tulid hommikust
peale end vee äärde seadma oma \VUgras{õngeritvadega} \textsuperscript{(45)}.

%%K t08-c2-09-2 p
\VUgras{Kala} \textsuperscript{(47)} haarab \VUgras{õngekonksu} kohe, kui nad kastavad
oma \VUgras{nööri} \textsuperscript{(46)} vette, ja vaevalt on neil aega \VUgras{sööta}
uuendada, kui uus ohver nööri otsas siputab.

%%K t08-c2-09-3 p
Siiski püüab \VUgras{võrgupüüdja} \textsuperscript{(48)}, kes viskab oma
\VUgras{paadist} \textsuperscript{(49)} \VUgras{heitevõrgu} keset \VUgras{jõge}
\textsuperscript{(50)}, rohkem kalu.

%%K t08-c2-10-1 p
10. --- Sügis on heade jalutuskäikude aeg: jalgsi, \VUgras{ratsa}
\textsuperscript{(52)}, \VUgras{jalgrattaga} \textsuperscript{(53)}, \VUgras{autoga}
\textsuperscript{(54)}. \VUgras{Teed} \textsuperscript{(51)} on puhtad ja kasutatakse ära
viimaseid ilusaid päevi, sest vaat et juba kogunevad \VUgras{pääsukesed}
\textsuperscript{(56)} katustele, et lennata kõigest jõust soojemate maade poole. Vana
\VUgras{pähklipuu} \textsuperscript{(57)} lehestik läheb kollaseks, \VUgras{pähklid}
kukuvad ja kõrged \VUgras{puud}, mis on rühmitatud nagu \VUgras{kimp}
\textsuperscript{(58)} \VUgras{kõrgendikul} \textsuperscript{(59)}, jäävad varsti
lehtedeta.

%%K t08-c2-11-1 p
11. --- \VUgras{Päike} \textsuperscript{(60)} tõuseb hiljem, päevad lühenevad, üsna
torkav külmake hakkab tunda andma \VUgras{hommikul}, ja vaat et juba lahkuvad
\VUgras{metskurvitsate} \textsuperscript{(61)} parved meie vähe külalislahketest
maadest.

\VUsaut{6.0mm}
\VUfilet{22mm}
